\section{Groups}\label{app:groups}
\textit{Groups} are used as a fundamental mathematical language to describe, understand and predict symmetries and physical phenomena within physics, particularly in quantum field theory. A group $G$  is a set of elements $\{g\}$ with an operation that assigns to every pair of elements a third element within that same set. The third element must satisfy \cite{Georgi:1999wka}
\begin{enumerate}
    \item Closure: If $f,g \in G$ then $h=fg\in G$.
    \item Associativity: For all $g,\,h,\,k,\in G$, $(gh)k=g(hk)$.
    \item Identity element: There exist an element, $e\in G$ such that for all $f\in G$ we have $ef=fe=f$.
    \item Inverse element: For every $f\in G$ there exist an inverse element such that $gg^{-1}=g^{-1}g=e$.
\end{enumerate}
Groups can be interpreted more visually as a multiplication table, as seen in Table \ref{tab:Z3} (for discrete group elements), this example is for the group $\mathrm{Z}_3$.
\begin{table}[h]
    \centering
    $
    \begin{array}{c||c|c|c} % Removed the last '|' here
        \backslash & e & a & b \\
        \hline\hline
        e & e & a & b \\
        \hline
        a & a & b & e \\
        \hline
        b & b & e & a  % Removed the \hline here
    \end{array}
    $
    \caption{Multiplication table of the finite group $\mathrm{Z}_3$.}
    \label{tab:Z3}
\end{table}
A group can be both finite and infinite. $\mathrm{Z}_3$ is a \textit{finite group} because it has a finite number of group elements, whereas an \textit{infinite group} will have an infinite number of group elements. The number of elements in a finite group is called the order of the group. Thus, $\mathrm{Z}_3$ is of order $3$ \cite{Georgi:1999wka}. Groups can also be categorised by being \textit{abelian} or \textit{non-abelian}. The operation of an abelian group has an extra axiom, restricting the multiplication law
\begin{enumerate}
    \item Commutative: For all $g,h\in G$, $gh=hg$.
\end{enumerate}
From the multiplication table for $Z_3$ \ref{tab:Z3} it can be seen that the group is abelian as well. 

\section{Representation} \label{app:repersentations}
A \textit{representation} of a group $G$ is a mapping, $D$ of the elements of the group onto a set of linear operators acting on a vector space $V$. This makes the abstract group easier to realise, hence representations are matrices (for finite groups). This is useful because without this a mathematical abstract object like a group, cannot directly be understood. However, through a representation of the group, we can realise it using one of the things we understand well: matrices \cite{Georgi:1999wka}. 

The group representations are useful because they live in linear space. This also means that we can always change the representation by performing a linear transformation. The mapping $D$ need to uphold two conditions \cite{Georgi:1999wka}
\begin{itemize}
    \item The representation of the identity element $e\in G$ must yield the identity operator in the respective vector space: $D(e)=1$. 
    \item The group's multiplication law is mirrored by the multiplication of their corresponding linear operators on the vector space: $D(g_1)D(g_2)=D(g_1g_2)$.
\end{itemize}
An example of this is the 3D rotation group: $\SO(3)$. The group has a representation that consist of $3\times 3$ matrices. The dimension of a representation is defined by the dimensions of the space that it acts upon \cite{Georgi:1999wka}. In this example the dimensions is thus $3$. Each group has multiple different representations and they don't have to have the same dimensions. A representation is regular if the dimensions of the representation is the order of the group \cite{Georgi:1999wka}. 

Representations can also be described by being \textit{reducible},\textit{ completely reducible} and \textit{irreducible}. To understand these definitions, we need to first understand invariant subspaces. An invariant subspace is when all mappings $D(g)$ on any vector $w$ in the subspace $W\subset V $ maps the vector back into the subspace $W$ \cite{Georgi:1999wka}
\begin{equation}
    D(g)w\in W\quad \forall\; g\in G, w\in W\;.
\end{equation}
A representation is reducible if it has a non-trivial invariant subspace. Non-trivial meaning that it has to be an invariant subspace which is not just $\{0\}$ or the entire set of $V$. An irreducible representation (irrep) is when the representation has no non-trivial invariant subspaces. Thus you cannot "break" the representation further down. A completely reducible representation is on block diagonal form 
\begin{align}
    \begin{pmatrix}
        D_1(g) & 0      & \dots  \\
        0      & D_2(g) & \dots  \\
        \vdots & \vdots & \ddots
    \end{pmatrix}
\end{align}
and can thus be written as a direct sum of irreducible representations
\begin{equation}
    D_1\oplus D_2 \oplus \cdots\;.
\end{equation}
An important property of finite groups is that their representations always is completely reducible \cite{Georgi:1999wka}.

Another feature of representations is that they can be unitary. A representation is unitary if all $D(g)$'s are unitary, where unitarity means that the operator upholds $\mathrm{O}^\dagger=\mathrm{O}^{-1}$. All representations of finite groups are equivalent to unitary representations \cite{Georgi:1999wka}.

\section{Subgroups}
A \textit{subgroup} $H$ is a group with a set of elements which all is a part of a group $G$. Then $H$ will be a subgroup of $G$. There are two trivial subgroups of $G$, namely the identity and the group itself. Subgroups thus capture a smaller symmetry structure inside a larger symmetry \cite{Georgi:1999wka}.

\section{Young Tableaux} \label{app:youngtab}
To enable us to index and construct irreducible representation of a group we utilise permutations and combinatorics through Young tableaus. We will start by considering it for the symmetry group, before utilising it for $\SU(N)$ groups. 

The goal is thus to index a simple representation of the symmetry group $\mathrm{S}_n$ which describes all possible permutation of $n$ different objects. Let us consider the symmetry group $\mathrm{S}_4$. This groups permutes the four numbers $1,2,3,4$. It is a non-abelian group. To consider their representation theory, we can describe it using combinatorics. To do this we can use conjugacy classes of the symmetric group which has a one-to-one correspondence with the simple representations. So we need to be able to label the conjugated classes. The definition is that conjugate classes consists of all permutations with the same cycle type. A cycle type is the length of a permutation. So for an example
\begin{align}
    (1)(234)
    \label{eq:cycle}
\end{align}
would leave the first number unchanged. The second number would be changed to the third, the third to the fourth and the fourth to the second. So $1234$ would change to $1342$. This is a combination of a 1-cycle and a 3-cycle. The conjugacy classes can be used to find the irreducible representations of the group. This can be done using Young tableaux \cite{Georgi:1999wka}
\begin{align}
    \ytableausetup{boxsize=0.5cm}
    \begin{ytableau}
        {} & {}\\
        {}\\
        {}
    \end{ytableau}
\end{align}
This is the same cycle as Eq. \ref{eq:cycle}. Each cycle type is represented in columns. Here we have the first row of length three and the second has of length one. This tableau thus represents the conjugacy class $(2,1^2)$, where the first denotes the boxes in first row, and that the two following rows have one box in each . Each tableau represents a different conjugacy class, another example could be the identity element $(4)$
\begin{align}
    \ytableausetup{boxsize=0.5cm}
    \begin{ytableau}
        {} & {}& {}&{}
    \end{ytableau}
\end{align}
which is four 1-cycles. The constructions of these diagrams are bounded by 
\begin{align}
    \lambda_1\ge\lambda_2\ge...\ge\lambda_n, \quad \sum_{i=1}^n\lambda_i=n\;,
\end{align}
for $\mathrm{S}_n$ where $\lambda_i$ denotes the row considered. Each tableau not only corresponds to a conjugacy class but also an irreducible representation \cite{YoungTableaux}. The dimension of the irreducible representation can be determined from the Young tableau,by dividing $n!$ with the product of the hook-lengths $h$ of all the boxes together 
\begin{align}
   f^\lambda=\frac{n!}{\prod_{(i,j)\in \lambda}h_ij} \;,
\end{align}
where the hook-length is calculated by adding $1$ (for the box it self) with all the boxes to the right and all the boxes below. So for $\mathrm{S}_4$ with $n!=4!$ the tableaus are
\begin{align}
    \ytableausetup{boxsize=0.5cm}
    &\begin{ytableau}
        {4} & {3}& {2}&{1}
    \end{ytableau}
    \quad\quad \quad
    \begin{ytableau}
        {4} & {2}&{1}\\
        {1}
    \end{ytableau}
    \quad\quad \quad\quad
    \begin{ytableau}
        {3} & {2}\\
        {2}&{1}
    \end{ytableau}
     \quad\quad \quad\quad\quad
    \begin{ytableau}
        {4} & {1}\\
        {2}\\
        {1}
    \end{ytableau}
    \quad\quad \quad\quad
    \begin{ytableau}
        {4} \\
        {3}\\
        {2}\\
        {1}
    \end{ytableau}
    \\
    &\quad\frac{4!}{4!}=1,\quad \quad\quad\frac{4!}{4\cdot 2\cdot 1\cdot 1}=3,\quad\frac{4!}{3\cdot 2\cdot 2\cdot 1}=2,\quad\frac{4!}{4\cdot 1\cdot 2\cdot 1}=3,\quad\frac{4!}{4!}=1,
\end{align}
where the numbers inside denotes the hook-length of that box \cite{westrich2011young}. 
Now if we instead of finite groups consider the continuos groups $\SU(N)$ the Young tableaux will prove to still be useful for finding the irreps of a specific group, we will come back to this after considering the basics of Lie groups. 

\section{Lie groups} \label{app:Lie}
\textit{Lie groups} serve as an important group within physics, because they describe continuous symmetries. Unlike the some of the discrete groups which we have discussed prior to this Lie groups allow for smooth and infinitesimal transformations. Considering Noether's theorem \cite{Peskin:1995ev} this continuous symmetries will result in physical conservation laws. Additionally, if the symmetries are promoted to become local gauge symmetries, they explain structure of the interactions of the theory such as the Standard Model and Grand Unified Theories \cite{Peskin:1995ev}. 
\subsection{Non-abelian group representations}
Lie groups are continuously generated groups, where "continuously" indicates that using infinitesimal elements of the group a general element can be reached. This also requires, that the group elements are arbitrarily close to the identity. An infinitesimal element can be written \cite{Georgi:1999wka}
\begin{equation}
    g(\alpha)=1+i\alpha^aT^a + \mathcal{O}^a\;,\label{eq:inftetrans}
\end{equation}
where $T^a$ is the \textit{generators} of the group and $\alpha^a$ is the infinitesimal group parameters. The generators span the space of the infinitesimal group transformations and they obey the commutation relation
\begin{equation}
    [T^a,T^b]=if^{abc}T^c\label{eq:comTa}
\end{equation}
which is a linear combination of generators. The vector space spanned by the generators, which obey commutation are called the \textit{Lie Algebra} \cite{Peskin:1995ev}. The structure constants $f^{abc}$ obey the \textit{Jacobi identity}
\begin{equation}
    f^{ade}f^{bcd}+f^{bde}f^{cad}+f^{cde}f^{abd}=0\;.\label{eq:jacobi}
\end{equation}
The commutation relations of the Lie algebra completely defines how the group element multiply close to the identity. This is the infinitesimal transformation. Thus, gauge theories, which only consider local transformations, the physics is fully determined by its algebra. However, two different Lie groups can share their Lie algebra, thus the commutation relations are shared but they have topological differences \cite{Peskin:1995ev}. When interested in gauge theories, these global difference can be ignored. 

To ensure that we have unitary representation, we consider \textit{compact} groups. This type of groups will have a finite number of generators \cite{Peskin:1995ev}.

As previously defined a group can be both non-Abelian and Abelian (see Section \ref{app:groups}). Thus if one of the generators of the group commutes with all other generators, it generates an independent continuous Abelian group \cite{Peskin:1995ev}. This type of group is called $\U(1)$ and describes phase rotations 
\begin{align}
    \psi\to e^{i\alpha}\psi\,.
\end{align}
If the algebra cannot be divided into mutually commuting subsets of generators, then the the Lie algebra is simple and belongs to a non-Abelian group. If it is possible to divide the generators in such subsets, then we have a semi-simple algebra, and can write the group structure as a product such as $\SU(4)\times\SU(2)\times\SU(2)$ \cite{Peskin:1995ev}.  

The so-called Cartan--Killing classification of simple, compact Lie algebras lead to the discovery of few possible groups which nearly all belong to the three infinite families of \textit{classical groups} \cite{Peskin:1995ev}
\begin{enumerate}
    \item \textbf{Unitary groups} $\SU(N)$ which are unitary transformations of complex $N-$dimensional vectors. It consists of $Ñ\times N$ unitary transformations with determinant one, and has $N\times N$ Hermitian matrices $t^a$ as generators. These have to satisfy $\tr t^a=0$. There are $N^2-1$ independent matrices satisfying this \cite{Peskin:1995ev}.
    \item \textbf{Orthogonal Groups} $\SO(N)$ also known as the rotation group in $N$ dimensions, this is a subgroup of unitary $N\times N$ transformations, which preserves the symmetric inner product. There are $N(N-1)/2$ generators. 
    \item \textbf{Symplectic Groups} $\mathrm{Sp}(N)$ this is a subgroup of unitary $N\times N$ transformations, which preserves the anti-symmetric inner product.  There are $N(N+1)/2$ generators. 
\end{enumerate}
Besides these there are the \textit{exceptional} Lie algebras $\mathrm{G}_2, \mathrm{F}_4, \E_6, \E_7,$ and $\E_8$. 

We will focus on the unitary groups, as these are the focus of the thesis. 

Let us consider the trace of the product of two generators
\begin{align}
    \tr T^a_R T^b_R\equiv D^{ab}\,.
\end{align}
where the subscript $R$ denotes the irrep. We choose a basis such that the $D^{ab}$ is proportional to the identity \cite{Peskin:1995ev}
\begin{equation}
    \tr{T^a_R T^b_R}=T(R)\delta^{ab}\;,\label{eq:Dynkinindex}
\end{equation}
where $T(R)$ is the Dynkin index \cite{srednicki2006qft}, and is a constant for each irrep $R$. By multiplying Eq. \eqref{eq:comTa} with another generator $T^c_R$ on the right and evaluating the trace, we obtain an expression for the structure constant, which we can write in terms of the normalisation from Eq. \eqref{eq:Dynkinindex}
\begin{equation}
    T(R)f_{abc}=-i\tr{[T^a_R ,T^b_R]T^C_R}\,,\label{eq:stuctureconstant}
\end{equation}
because of trace cyclicity
\begin{align}
    \tr{[T^a_R ,T^b_R]T^C_R}=\tr(T^a_RT^b_RT^C_R-T^b_RT^a_RT^C_R)=\tr(T^b_RT^c_RT^a_R-T^c_RT^b_RT^a_R)=\tr{[T^b_R ,T^c_R]T^a_R}\,,
\end{align}
permutations of the indices of structure constant does not change the sign. However, interchanging indices does (due to the nature of the commutator, evident from \eqref{eq:comTa}). Thus, the structure constant is completely anti-symmetric \cite{Georgi:1999wka}.

For each irrep of the group, there will be a conjugate representation. This will be noted with a bar: $\overline R$. Using the infinitesimal transformation defined in Eq. \eqref{eq:inftetrans} we can write the complex conjugate 
\begin{align}
    \psi\to(1+i\alpha^a(T^a_R))\psi\\ \psi^*\to(1+i\alpha^a(T^a_R)^*)\psi^*\,.
\end{align}
Thus we can write the conjugate representation as $T^a_{\overline R}=-(T^a_R)^T$. This ensures that combinations of a field with its complex conjugates transforms as singlet under the gauge group \cite{Peskin:1995ev}.

The unitary group has a so-called \textit{fundamental} representation. This will be the $N-$dimensional complex vector. Hence, for $N\ge2$ this irrep will be complex, there is a corresponding independent anti-fundamental representation, which will be the $\overline N-$dimensional complex vector. 

Another type of representation which is irreducible is the adjoint representation of the group. This representation is generated from the structure constants \cite{Georgi:1999wka}
\begin{align}
    (T^a_\text{adj})_{bc}=-if_{abc}\,.
\end{align}  
This can be seen from considering the condition on the generator to satisfy the Lie algebra in Eq. \eqref{eq:comTa}. For the adjoint representation, this is just a rewriting of the Jacobi identity in Eq. \eqref{eq:jacobi} \cite{Georgi:1999wka}. Thus, the Jacobi identity proves that the adjoint representation is a valid representation. Additionally, as an artefact of the structure constants being completely anti-symmetric the representation is real. The dimensions of this representation for the unitary group is $d(\text{adj})=N^2-1$ \cite{Peskin:1995ev}.

A characteristic of the representations besides the Dynkin index is the  \textit{quadratic Casimir operator} $C_2(R)$. This factor arises when considering the operator
\begin{align}
    T^2=T^a T^a\,,
\end{align}
which commute with all generators
\begin{align}
    [T^2,T^b]=0\,.
\end{align}
This means that it is an invariant of the algebra and thus a constant for each irrep. Thus, we can write the operator as 
\begin{align}
    T^a_R T^a_R=C_2(R)\mathbb{I}\,,
\end{align}
with the dimensions of the unit matrix being the dimensions of the irrep. The two constants of the irreps are related as \cite{Peskin:1995ev}
\begin{align}
    C_2(R)=\frac{d(\text{adj})}{d(R)}T(R)\,.
\end{align}
Using the convention that 
\begin{equation}
    \tr{T^a_N T^b_N}=\frac{1}{2}\delta^{ab}\;.
    \label{eq: Dynkin index fund. rep}
\end{equation}
we can write the values of these constants for $\SU(N)$. An overview can be seen in Table \ref{tab:SUN}
\begin{table}[h]
    \centering
    \renewcommand{\arraystretch}{2.2} % Prevents fractions from looking squished
    \ytableausetup{boxsize=0.25cm, centertableaux} % Set globally once to clean up the rows
    $
    \begin{array}{c||c|c|c} % Changed lc|| to c|| to remove the ghost column
        \text{Representation} & d(R) & T(R) & C_2(R) \\
        \hline\hline
        \begin{ytableau}{}\end{ytableau} \,\;/\,\; \overline{\begin{ytableau}{}\end{ytableau}\rule{0pt}{7.5pt}} & N & \dfrac{1}{2} & \dfrac{N^2 - 1}{2N} \\
        \hline
        \mathbf{adj} & N^2 - 1 & N & N \\
        \hline
        \begin{ytableau}{}&{}\end{ytableau} & \dfrac{N(N+1)}{2} & \dfrac{N+2}{2} & \dfrac{(N-1)(N+2)}{N} \\
        \hline
        \begin{ytableau}{} \\ {}\end{ytableau} & \dfrac{(N-1)N}{2} & \dfrac{N-2}{2} & \dfrac{(N-2)(N+1)}{N} % Fixed to vertical boxes for anti-symmetric
    \end{array}
    $
    \caption{Dimensions $d(R)$, Dynkin indices $T(R)$, and quadratic Casimir invariants $C_2(R)$ for the fundamental, adjoint, two-index symmetric, and two-index anti-symmetric representations.}
    \label{tab:SUN}
\end{table}

If we now consider anti-commutator relations in place of the commutator in Eq. \eqref{eq:stuctureconstant} we have another invariant symbol 
\begin{equation}
    A(R)d^{abc}\equiv \frac 12 \tr{T^a_R\{T^b_R,T^c_R\}}\;,
\end{equation}
with $A(R)$ being the anomaly coefficient \cite{srednicki2006qft} and $d^{abc}$ being a symmetric tensor. Using the cyclic property of the trace, we have
\begin{equation}
    A(R)=-A(\bar R)\;.
    \label{eq:anomaly coefficient}
\end{equation}
For real or pseudoreal representation the anomaly coefficient is thus $A(R)=0$. Thus, for chiral theories it is non-trivial for this coefficient to vanish, ensuring an anomaly free gauge sector.


\subsection{Young tableaux for SU(N)} \label{app:youngtabSUN}
As mentioned in Section \ref{app:youngtab} the Young tableau approach can be generalised to be used to find irreps of a Lie Algebra. This is done by considering the highest weight vector, $\Lambda$ \cite{Georgi:1999wka}. This weight vector is defined from the fundamental weights $\{\eta^k\}^r_{k=1}$, which is related to the root space basis vectors \footnote{The vector-space basis of the root space is written in terms of the simple roots $\{\alpha_i\}^r_{i=1}$. This root space comes from the subalgebra of a Lie algebra: The Cartan subalgebra. This vector-space basis together with the alternative basis $\{\eta^k\}^r_{k=1}$, showing that the fundamental weights are dual with the simple roots \cite{YoungTableaux}.}
\begin{align}
    \Lambda\equiv m_1\eta^1 +...+m_r\eta^r\;,
\end{align} 
where $(m_1,..,m_r)$ are the Dynkin integers. These integers is what we use to distinguish between different irreps. Each irrep will have its own distinct values of $(m_1,..,m_r)$, where the Dynkin integers must be non-negative. To understand the irreps of a certain Lie Algebra, we can now use Young tableaux, to determine the dimensions of these irreps and the decomposition of a direct product of two irreps. We can do this by considering the value of the $m_j$'s as the number of columns of length $j$. For one single box $\begin{ytableau}{} \end{ytableau}$ we will have $m_1=1$ and $m_2,...,m_r=0$. This can also be written as $(1,0,0,...,0)$. For $(0,2,0,0,...,0)$ we would have two columns of two boxes each \cite{YoungTableaux}
\begin{align}
    \ytableausetup{boxsize=0.5cm}
    \begin{ytableau}
        {} & {}\\ {}&{}
    \end{ytableau}
\end{align}
The number of rows is restricted by the group
\begin{align}
    \text{\# of boxes} =\sum^r_{j=1}j m_j\;,
\end{align}
for $SU(N)$, $r=N-1$, this means that for $SU(4)$ e.g. the maximum number of boxes will be 3 in each column for non-trivial representations. This is because a column of $n$ boxes would correspond to the determinant representation which is simply 1 \cite{YoungTableaux}. $SU(N)$ has two invariant tensors $\delta^k_l$ and $e^{abc...n}$, this is different from other Lie algebras, and therefore the Young tableaux scheme is mostly used for $SU(N)$ \cite{YoungTableaux}. The rows describe symmetric indices, and the columns describe the antisymmetric indices.   

Now to compute the dimensions of the irrep, the hook length $h_{ij}$ is used \cite{YoungTableaux}
\begin{align}
    \dim(\Lambda)=\prod_{i,j\in \Lambda}\frac{N+j-i}{h_ij}\;,
\end{align}
where $j$ denotes column and $i$ denotes row. Using this formula, we consider an example, let us focus on the Young tableau $(1,1,0,0,...,0)$
\begin{align}
    \ytableausetup{boxsize=0.5cm}
    \dim\left(\begin{ytableau}
        {} & {}\\ {}
    \end{ytableau}\right)=\begin{ytableau}
        {N} & {M}\\ {P}
    \end{ytableau}\; \text{ divided by } \;\begin{ytableau}
        {3} & {1}\\ {1}
    \end{ytableau}=\frac{1}{3}N(N^2-1)\;,
\end{align}
where $M=N+2-1=N+1$ and $P=N+1-2=N-1$. Now that we can compute the dimensions of a specific irrep from the Young tableau, the next step is to use the Young tableaux to compute direct products. Understanding how to decompose two or more representations into a direct sum of irreps, is best shown by an example. Let us consider the direct product of $(0,0,1,0,...,0)$ and $(1,1,0,0,...,0)$
\begin{align}
    \ytableausetup{boxsize=0.5cm}
    \begin{ytableau}
        {} \\ {}\\{}
    \end{ytableau}
    \otimes \begin{ytableau}
        {a} &{a}\\ {b}
    \end{ytableau}\;,
\end{align}    
the first step is to write up the two Young tableaux for the representations in the direct product. In the second Young tableau one writes a new letter in each row, so $a,b$ in the first two, in the third row $c$, and so on. For a row with multiple boxes, all boxes in the row gets the same letter. Next, the first box (the box furthest to the top and to the right) is attached to the three boxes of the first tableau, in all possible combinations \cite{YoungTableaux}
 \begin{align}
    \ytableausetup{boxsize=0.5cm}
    \begin{ytableau}
        {} \\ {}\\{}
    \end{ytableau}
    \otimes \begin{ytableau}
        {a} &{a}\\ {b}
    \end{ytableau}=\left(
        \begin{ytableau}
        {} &{a}\\ {}\\ {}
        \end{ytableau}\oplus 
        \begin{ytableau}
        {} \\ {}\\ {}\\{a}
        \end{ytableau} 
    \right)
    \otimes \begin{ytableau}
        {a}\\ {b}
    \end{ytableau}\;.
\end{align} 
After this the next box can be attached,
 \begin{align}
    \ytableausetup{boxsize=0.5cm}
    \begin{ytableau}
        {} \\ {}\\{}
    \end{ytableau}
    \otimes \begin{ytableau}
        {a} &{a}\\ {b}
    \end{ytableau}=\left(
        \begin{ytableau}
        {} &{a}&{a}\\ {}\\ {}
        \end{ytableau}\oplus
        \begin{ytableau}
        {} &{a}\\ {}\\ {}\\{a}
        \end{ytableau}\oplus  
        \begin{ytableau}
        {} &{a}\\ {}\\ {}\\{a}
        \end{ytableau}\oplus
        \begin{ytableau}
        {} \\ {}\\ {}\\{a}\\{a}
        \end{ytableau} 
    \right)
    \otimes \begin{ytableau}
        {b}
    \end{ytableau}\;.
\end{align} 
however not all of these tableaux is allowed. Having the same letter in the same column is prohibited, thus rendering the last tableau wrong \cite{YoungTableaux}. The two middle tableaux is identical \footnote{This is where the notation of $a,b$ is important. If two tableaux have the same shape, but different letters in different boxes. They would both be kept. }, and we will only keep one, thus the expression reduces to
 \begin{align}
    \ytableausetup{boxsize=0.5cm}
    \begin{ytableau}
        {} \\ {}\\{}
    \end{ytableau}
    \otimes \begin{ytableau}
        {a} &{a}\\ {b}
    \end{ytableau}=\left(
        \begin{ytableau}
        {} &{a}&{a}\\ {}\\ {}
        \end{ytableau}\oplus
        \begin{ytableau}
        {} &{a}\\ {}\\ {}\\{a}
        \end{ytableau}
    \right)
    \otimes \begin{ytableau}
        {b}
    \end{ytableau}\;.
\end{align}
Continuing, with the last box
 \begin{align}
    \ytableausetup{boxsize=0.5cm}
    \begin{ytableau}
        {} \\ {}\\{}
    \end{ytableau}
    \otimes \begin{ytableau}
        {a} &{a}\\ {b}
    \end{ytableau}=
        \begin{ytableau}
        {} &{a}&{a}&{b}\\ {}\\ {}
        \end{ytableau}\oplus
        \begin{ytableau}
        {} &{a}&{a}\\ {}&{b}\\ {}
        \end{ytableau}\oplus
        \begin{ytableau}
        {} &{a}&{a}\\ {}\\ {}\\{b}
        \end{ytableau}\oplus
        \begin{ytableau}
        {} &{a}&{b}\\ {}\\ {}\\{a}
        \end{ytableau}\oplus
        \begin{ytableau}
        {} &{a}\\ {}&{b}\\ {}\\{a}
        \end{ytableau}\oplus
        \begin{ytableau}
        {} &{a}\\ {}\\ {}\\{a}\\{b}
      \end{ytableau}\;,
\end{align}
once again not all diagrams are allowed. The additional rule is that, when adding a box, one needs to make sure that from the right, the letter from a higher row number does not precede that of a lower row number \cite{YoungTableaux}. So in this case of $a$'s and $b$'s, we need to count the number of boxes with each letter, staring from the first row at the right, moving to the left and doing the same in the following rows. If at any point in this counting, there are more $b$'s than $a$'a, the tableau should be removed. This is the same for more $c$'s than $b$'s and so on. This also means that the first and the fourth tableau should be deleted  
\begin{align}
\begin{ytableau}
        {} \\ {}\\{}
    \end{ytableau}
    \otimes \begin{ytableau}
        {a} &{a}\\ {b}
    \end{ytableau}=
    \ytableausetup{boxsize=0.5cm}
        \begin{ytableau}
        {} &{a}&{a}\\ {}&{b}\\ {}
        \end{ytableau}\oplus
        \begin{ytableau}
        {} &{a}&{a}\\ {}\\ {}\\{b}
        \end{ytableau}\oplus
        \begin{ytableau}
        {} &{a}\\ {}&{b}\\ {}\\{a}
        \end{ytableau}\oplus
        \begin{ytableau}
        {} &{a}\\ {}\\ {}\\{a}\\{b}
      \end{ytableau}\;,
\end{align}
one continues attaching the boxes of the second tableau in the direct product onto the first, until none remains. Thus, in this example we end up with four irreps, and to check that the dimensions fit, we can calculate the dimensions of the irreps. To do this, we need to assume a group, let us take $\SU(5)$. This also, changes on of the diagrams. As mentioned previously the maximum length of a column should be $N-1$. Thus, the last diagrams first column would be invariant under $\SU(5)$, and we therefore delete it 
\begin{align}
\begin{ytableau}
        {} \\ {}\\{}
    \end{ytableau}
    \otimes \begin{ytableau}
        {} &{}\\ {}
    \end{ytableau}=
    \ytableausetup{boxsize=0.5cm}
        \begin{ytableau}
        {} &{}&{}\\ {}&{}\\ {}
        \end{ytableau}\oplus
        \begin{ytableau}
        {} &{}&{}\\ {}\\ {}\\{}
        \end{ytableau}\oplus
        \begin{ytableau}
        {} &{}\\ {}&{}\\ {}\\{}
        \end{ytableau}\oplus
        \begin{ytableau}
        {} 
      \end{ytableau}\;,
\end{align}
with the dimensions
\begin{align}
    10 \times 40 = 280+70 + 45+5 \;,
\end{align} 
where on both sides the dimension is equal to $400$, thus it is consistent. Using the Dynkin integer notation the direct product can be written as
\begin{align}
    (0,0,1) \otimes (1,2) = (1,1,1)\oplus (2,0,0,1) \oplus (0,1,0,1)\oplus (1) \;.
\end{align} 
This method is useful within physics, because it is a systematic way to decompose direct products of representations. This is especially useful GUTs, where it is needed to decompose large representations into the representations of the Standard Model.

